\section{Chapter 2. Methodology}\label{chapter-2.-methodology}

\subsection{2.1. Data sources}\label{data-sources}

\subsubsection{2.1.1. Air quality: OpenAQ}\label{air-quality-openaq}

The primary observation source for this study is the OpenAQ platform \cite{openaq_platform}, an open-access repository that aggregates air quality measurements from government reference monitors, research-grade instruments, and low-cost sensor (LCS) networks worldwide. Data are accessed through the OpenAQ v3 REST API at the base endpoint \texttt{https://api.openaq.org/v3}, with requests authenticated by a personal API key supplied in the \texttt{X-API-Key} HTTP header. The spatial query specifies a 25 km radius centred on Almaty (43.2389°N, 76.8897°E), which encompasses the full built-up area of the city including the foothill districts to the south and the industrial and residential zones to the north and west. Within this radius, 192 station records were registered in the OpenAQ index as of the day-2 data snapshot (2026-04-20); of these, 190 were actively reporting PM2.5 during the study window. The provider mix at that snapshot was 153 AirGradient devices, 38 Clarity Node-S units, and one AirNow federal reference station operated by the United States Embassy. All 192 instruments are classified as low-cost sensors under the taxonomy introduced in §1.3.

The OpenAQ v3 API paginates location queries in batches of up to 100 locations per request, requiring two consecutive calls to retrieve all 192 stations in the spatial query. Each per-sensor measurement request returns up to 100 hourly records per call. For a sensor with 562 days of continuous data, approximately six API calls are required to retrieve the full measurement archive; in practice, most AirGradient sensors came online in early 2025, so the typical archive spans 300--400 days and requires four to five calls per sensor. The ingest module (\texttt{src/ingest/openaq.py}) implements an exponential back-off retry strategy with a base wait of 5 seconds, a backoff factor of two, and a maximum of eight retry attempts per call. Under this schedule, a failed request can wait up to 5 × (2$^8$ − 1) = 1,275 seconds before the final retry; in practice, the API returns HTTP 429 (rate limit) or 503 (temporary unavailability) errors on fewer than 2\% of calls. With 190 sensors each requiring four to five API calls, the full ingest run takes approximately three hours on a residential internet connection.

The ingest design is idempotent. Each sensor's measurement archive is written to a dedicated Parquet file at \texttt{data/raw/pm25/\{sensor\_id\}.parquet}. At the start of each run, the ingest script checks for the presence of this file; if the file exists, the sensor is skipped unless the \texttt{-\/-force} flag is passed at the command line. This skip-if-present logic ensures that a failed run can be resumed without duplicating already-retrieved data and without incurring unnecessary API quota consumption. The force-refresh mode overwrites the existing file in place, so the raw tier contains at most one file per sensor at any time.

The OpenAQ platform publishes all data under the Creative Commons Attribution 4.0 International licence (CC BY 4.0). This licence permits unrestricted use, redistribution, and the creation of derivative works, including the trained forecasting models and the Streamlit application produced in this study, provided that the source is attributed. No additional data-sharing agreement or ethical approval is required to access OpenAQ data, as the platform aggregates measurements that sensor operators have already consented to share publicly.

\subsubsection{2.1.2. Meteorology: Open-Meteo historical archive}\label{meteorology-open-meteo-historical-archive}

Meteorological reanalysis data are retrieved from the Open-Meteo historical weather archive \cite{zippenfenig_openmeteo_2023} at the endpoint \texttt{https://archive-api.open-meteo.com/v1/archive}. The archive derives from the ECMWF ERA5 global reanalysis and provides hourly gridded output interpolated to any user-specified point. No API key is required for non-commercial use; the public rate limit permits 10,000 requests per day, which is sufficient for a single point query covering the full 562-day study window in one call. The query is issued for Almaty centre (43.2389°N, 76.8897°E) and returns 13,488 hourly rows spanning 2024-10-01 through 2026-04-15.

Seven meteorological variables are retrieved. The first is \texttt{temperature\_2m}: screen-level air temperature at 2 m above ground level (AGL), reported in degrees Celsius. Temperature at screen level is the standard meteorological near-surface temperature measurement; it determines the thermal stratification of the surface boundary layer and is the primary driver of nocturnal inversion formation when radiative cooling drives the surface temperature below the air temperature aloft. The second variable is \texttt{relative\_humidity\_2m}: relative humidity at 2 m AGL, reported as a percentage. Relative humidity is relevant to PM2.5 measurement by light-scattering LCS instruments because hygroscopic growth of aerosol particles at high humidity causes the instrument to over-report mass concentration relative to the gravimetric standard. While this study does not apply a humidity correction to the raw LCS data, relative humidity is included as a predictor so that the model can learn the association between high humidity and inflated sensor readings, partially compensating for the uncorrected hygroscopic effect.

The third variable is \texttt{pressure\_msl}: mean sea-level pressure in hectopascals. Synoptic-scale pressure patterns modulate the frequency and intensity of blocking anticyclones over Central Asia, which are the meteorological precondition for the multi-day inversion episodes responsible for the highest PM2.5 concentrations in Almaty. Pressure at mean sea level serves as a proxy for the prevailing synoptic regime without the altitude correction complications that arise when using station pressure at Almaty's 847 m elevation. The fourth and fifth variables are \texttt{wind\_speed\_10m} and \texttt{wind\_direction\_10m}: wind speed in m/s and direction in degrees clockwise from north, both at 10 m AGL. Winds at 10 m represent the near-surface flow within the roughness sublayer and the lower surface layer; they are weaker than the gradient-level flow typically observed at 500--1,000 m AGL. For Almaty, the critical flow distinction is between northerly and north-westerly winds (which advect cleaner steppe air into the city) and south-easterly flows associated with valley recirculation and thermal drainage from the Zailiiskiy Alatau foothills. The sixth variable is \texttt{precipitation}: hourly total precipitation in millimetres. Precipitation, including both rain and snow, removes PM2.5 from the atmosphere through wet deposition (both below-cloud scavenging and in-cloud incorporation). Even modest precipitation of 0.5 mm or more in a given hour is associated with measurable reductions in PM2.5 in subsequent hours, making precipitation a useful binary and continuous predictor of inversion-break events.

The seventh variable is \texttt{boundary\_layer\_height} (BLH): the depth of the planetary boundary layer from the surface to the mixing height, in metres. The BLH diagnostic in the ERA5-based archive is a modelled output derived from the ECMWF Integrated Forecasting System, not a direct measurement. The ERA5 BLH is diagnosed using the bulk Richardson number criterion: the top of the boundary layer is defined as the height at which the bulk Richardson number, computed from the surface upward, exceeds a threshold of 0.25. The resulting BLH values in the training window range from approximately 10 m during deep winter inversions to 3,935 m during convective summer afternoons. This range of nearly four orders of magnitude motivates the log-transform feature described in §2.3.3. It is recognised in the ERA5 literature that the Richardson-number-based BLH diagnostic tends to underestimate the boundary-layer depth during stable stratification and to show higher uncertainty during nocturnal low-level jet conditions \cite{guo2021blh}. For this reason, the present study incorporates BLH as one of 27 predictors rather than as the sole meteorological predictor; its influence on model output is evaluated in the feature importance analysis presented in §3.4.

\subsubsection{2.1.3. Baseline forecast: CAMS via Open-Meteo air-quality API}\label{baseline-forecast-cams-via-open-meteo-air-quality-api}

The Copernicus Atmosphere Monitoring Service (CAMS) Global Atmospheric Composition Forecast provides the third external baseline for model comparison. CAMS-Global is a global chemistry-transport model operated by ECMWF on behalf of the European Union's Copernicus Earth Observation Programme. The model runs at approximately 40 km horizontal resolution on a 0.4° × 0.4° latitude-longitude grid and uses 137 vertical levels from the surface to 0.01 hPa. The aerosol scheme assimilates satellite-retrieved aerosol optical depth (AOD) from MODIS Terra and Aqua instruments (at 0.55 μm) and from the Sentinel-5P TROPOMI instrument, but does not assimilate ground-level LCS observations of the kind used in this study. CAMS PM2.5 at the surface is estimated from the model's total column of particulate matter, which aggregates dust, sea salt, organic carbon, black carbon, and sulphate aerosol species.

The operational CAMS forecast is initialised twice daily at 00 UTC and 12 UTC and covers a 120-hour (five-day) forecast range. For the historical backtesting period used in this study (2024-10-01 through 2026-04-15), however, the live forecast output is not available in a consistent time series. The present study therefore accesses the CAMS reanalysis product through the Open-Meteo air-quality historical API at the endpoint \texttt{https://air-quality-api.open-meteo.com/v1/air-quality} with the parameter \texttt{domains=cams\_global}. This endpoint returns the CAMS analysis and hindcast product rather than the prospective forecast, ensuring temporal consistency with the ERA5-based meteorological archive retrieved from the same Open-Meteo platform. The CAMS PM2.5 variable is retrieved at hourly resolution for the full 13,488-hour study window; the hourly output is collocated to the Almaty centre grid point (approximately 43.0°N, 76.5°E, the nearest CAMS model grid point to the city centre given the 0.4° grid spacing). The \textasciitilde30 km offset between this grid point and the true city centre (43.2389°N, 76.8897°E) is a consequence of the model's horizontal resolution and may introduce a small systematic bias in CAMS forecasts during local inversion events that are spatially confined to the Almaty basin. The CAMS data are published under CC BY 4.0.

The role of CAMS in the evaluation framework is as an evidence-based external baseline: it represents the best openly available PM2.5 forecast for a user who lacks access to local sensor data. Its inclusion in the comparison allows the present study to quantify the marginal value added by training a local data-driven model on the 190-sensor LCS network, rather than relying on the global CTM output. The expectation, stated as hypothesis H2 in §1.5, is that CAMS will fail to capture sub-grid inversion dynamics and will therefore show negative R² on the test set.

\subsubsection{2.1.4. Training window}\label{training-window}

The study window spans 2024-10-01 through 2026-04-15, a total of 562 days corresponding to 13,488 hourly meteorological rows. The temporal boundaries were determined by two constraints. The upper boundary (April 2026) was set to be as late as possible before the thesis writing period, maximising the amount of winter-season data available for training and evaluation. The lower boundary (October 2024) was forced by data availability: of the 192 OpenAQ sensors in the spatial domain, only one has continuous data extending back to 2023-24. The bulk of the AirGradient network was deployed between January and May 2025, and the sensor count remained below 50 stations before October 2024. Training on the period before October 2024 would have required using a much smaller and potentially less representative spatial sample, reducing the reliability of the city-median target variable.

The October 2024 start date carries a secondary benefit. The AirGradient network was most spatially dense and operationally stable during the 2024-25 and 2025-26 heating seasons, with 153 active AirGradient devices contributing to the median computation for the majority of winter hours. Sensor dropouts, periods where a device goes offline due to firmware updates, power failures, or connectivity issues, are less consequential when the median is computed over a large pool of instruments, because individual sensor absences do not shift the cross-sectional median. The training window thus encompasses two full heating seasons (October 2024 -- March 2025 and October 2025 -- April 2026) and one summer period (April -- September 2025), providing the model with examples of the full annual PM2.5 cycle.

The restriction to a single city and a 19-month window limits the generalisability of the trained models. XGBoost models trained on Almaty data cannot be expected to transfer directly to other Central Asian cities such as Bishkek or Tashkent, where topography, fuel mix, and sensor network density differ substantially. Within Almaty, the models are trained on city-median concentrations rather than station-level data, which means they are not designed to resolve within-city spatial gradients. These scope constraints are acknowledged in §4.4.

\subsubsection{2.1.5. Data storage architecture}\label{data-storage-architecture}

The pipeline uses a two-tier storage design that separates immutable raw data from derived analytical products. The design principle is that the raw tier can always be regenerated from the source APIs, while the derived tier can always be regenerated from the raw tier; no transformation result is treated as irreplaceable.

The first tier, the raw tier, stores per-sensor Parquet files for PM2.5 measurements and single-window Parquet files for meteorology and CAMS. Each per-sensor PM2.5 file resides at \texttt{data/raw/pm25/\{sensor\_id\}.parquet} and uses a fixed three-column schema: \texttt{timestamp} (UTC, stored as an integer Unix timestamp in nanoseconds by the Pandas Parquet writer), \texttt{sensor\_id} (string, matching the OpenAQ v3 location identifier), and \texttt{value} (float64, PM2.5 concentration in μg/m³). The Parquet files are encoded with Snappy compression, which achieves typical compression ratios of 3:1 to 5:1 for floating-point time-series data and decompresses in negligible time on modern hardware. The meteorological archive is stored at \texttt{data/raw/meteo.parquet} (13,488 rows × eight columns including timestamp) and the CAMS archive at \texttt{data/raw/cams.parquet} (13,488 rows × two columns). All three Parquet objects are listed in \texttt{.gitignore} and are not committed to version control.

The second tier, the derived tier, is a SQLite database at \texttt{data/almaty\_aq.db} containing three tables. The \texttt{pm25} table holds the per-sensor raw readings after parsing from Parquet (columns: \texttt{sensor\_id} TEXT, \texttt{timestamp} TEXT ISO-8601, \texttt{value} REAL). The \texttt{meteo} table holds the seven meteorological variables plus timestamp (eight columns, 13,488 rows). The \texttt{cams} table holds the CAMS PM2.5 values plus timestamp (two columns, 13,488 rows). The conversion from Parquet to SQLite is performed by \texttt{src/db.py}, which is idempotent: it reads the Parquet files, drops and recreates each table on each run, and inserts the full dataset. The SQLite file is also gitignored.

The SQLite layer serves as the interface for all downstream processing. Feature engineering (\texttt{src/features/pipeline.py}) queries the three tables via SQL joins and returns a Pandas DataFrame; model training and the Streamlit application then operate on this DataFrame. SQLite was chosen over a client-server database for two reasons. First, the absence of a server process makes the pipeline fully portable: it runs identically on a development laptop, in a Google Colab notebook, and on Streamlit Community Cloud, which does not support long-running background processes or network-exposed services. Second, SQLite's write-ahead log (WAL) mode, enabled by \texttt{PRAGMA\ journal\_mode=WAL}, provides concurrent read access, which allows the Streamlit application to query the database for live forecast display while a background data-refresh script is simultaneously writing new rows. For the data volumes in this study, 578,993 raw sensor-hour records loading into three tables totalling under 50 MB, SQLite provides sufficient read and write throughput without the operational overhead of PostgreSQL or any other client-server system.

The Parquet format for the raw tier was chosen for its columnar storage layout. Columnar encoding enables fast aggregation queries, for example, computing the cross-sectional median across 190 sensor columns for a single timestamp requires reading only the relevant column slices rather than the full row-oriented file. In Pandas, the \texttt{df.groupby(\textquotesingle{}timestamp\textquotesingle{}).median()} operation on a 578,993-row Parquet-derived DataFrame executes as a single vectorised operation without loading any non-value columns. The net effect is that a fresh full rebuild of the database from API downloads takes approximately three hours of network time (dominated by OpenAQ rate limits) but less than 90 seconds of local computation time. This separation of network-bound and compute-bound phases supports reproducibility: any user with valid API credentials and the repository code can rebuild the full dataset and retrain all models from scratch.

\subsection{2.2. System architecture}\label{system-architecture}

The pipeline is organised into five sequential stages: ingest, store, feature engineering, modelling, and application. Each stage communicates with the next through a well-defined interface (Parquet files or SQLite tables), and no stage imports directly from a non-adjacent stage. This separation of concerns prevents the ad hoc shortcuts, reading raw Parquet directly in the model training script, or querying the database from the ingest module, that create hidden dependencies and make the pipeline difficult to audit or extend.

\begin{verbatim}
+-----------------------+      +-------------------------+
|  OpenAQ v3  (LCS)     |      |  Open-Meteo archive     |
|  PM2.5 hourly         |      |  Meteo hourly + BLH     |
+----------+------------+      +------------+------------+
           |                                |
           |  src/ingest/openaq.py          |  src/ingest/meteo.py
           |  (locations + per-sensor       |  (single call per window)
           |   measurements, idempotent)    |
           v                                v
+---------------------------------------------------------+
|                data/raw/*.parquet                        |
|  (immutable snapshots; gitignored; re-creatable)         |
+------------------------+--------------------------------+
                         |  src/db.py
                         v
                +------------------+       +---------------------+
                | data/almaty_aq   | <---- |  Open-Meteo CAMS    |
                |  .db  (SQLite)   |       |  src/ingest/cams.py |
                +--------+---------+       +---------------------+
                         |  src/features/pipeline.py
                         v
                +------------------+
                | feature matrix   |  ---> src/models/baselines.py
                | (hourly, tidy)   |        src/models/xgboost_model.py
                +--------+---------+        src/models/evaluate.py
                         |
                         v
                +------------------+
                | src/app/main.py  |  (Streamlit: map + forecast card)
                +------------------+
\end{verbatim}

The ingest idempotency design merits further technical description. The OpenAQ ingest module queries the list of all locations in the spatial domain at the start of each run and writes the result to \texttt{data/raw/locations.parquet}. For each sensor in this list, the module checks whether a per-sensor Parquet file already exists at the expected path. If the file is present, the sensor is skipped entirely, no API call is made, no data is read. If the file is absent, the module fetches the full measurement archive for that sensor in paginated batches of 100 records per call and writes the concatenated result to the Parquet path. The \texttt{-\/-force} flag bypasses the existence check and overwrites all sensor files. This design ensures that the common incremental case (adding newly deployed sensors to an existing dataset) requires only one API call per new sensor rather than re-fetching the full 190-sensor archive. For the meteorological and CAMS archives, the idempotency model is simpler: these single-window files are always overwritten on each run, because they are requested as a single API call and the cost of re-fetching is trivial compared to the per-sensor PM2.5 archive.

The SQLite WAL-mode architecture is worth noting in the context of the Streamlit deployment. Streamlit Community Cloud runs the application server and any background scripts within a single virtual machine. When the application is running and a user is viewing the forecast dashboard, the Streamlit server process holds an open read connection to \texttt{almaty\_aq.db}. If a data-refresh script is executed concurrently to ingest new sensor data and update the database, WAL mode allows the reader to continue with a consistent snapshot of the database state at the moment the read transaction opened, while the writer proceeds independently. Without WAL mode, the SQLite default exclusive-write lock would block the reader for the duration of the write, causing the application to return a database-locked error. The WAL approach eliminates this failure mode at no cost in data consistency.

\subsection{2.3. Preprocessing and feature engineering}\label{preprocessing-and-feature-engineering}

\subsubsection{2.3.1. Station filtering and city-median aggregation}\label{station-filtering-and-city-median-aggregation}

Raw PM2.5 observations arrive from 190 active sensors across 192 registered locations in the OpenAQ spatial query. Before constructing the modelling target, two sequential filters are applied to the raw readings. The first filter discards all readings outside the physically plausible interval {[}0, 500{]} μg/m³. Readings below zero are physically impossible, a negative mass concentration has no physical interpretation, and in practice represent firmware overflow events or calibration errors affecting fewer than 0.01\% of records. Readings above 500 μg/m³ account for fewer than 0.04\% of the archive and are attributable to laser-scattering sensor saturation caused by contamination of the optical sensing chamber; they do not represent true atmospheric PM2.5 concentrations. The 500 μg/m³ threshold was chosen to be conservative: the WHO 24-hour guideline is 15 μg/m³ \cite{who_aqg_2021}, and even during Almaty's most severe documented inversion episodes, independent reference measurements have not exceeded 400 μg/m³.

The second filter requires a minimum of three sensors to report a valid reading in a given clock hour. Hours where fewer than three independent sensors agree produce a city-median that is vulnerable to individual station idiosyncrasies. A single malfunctioning device or a station located immediately downwind of a local emission source (vehicle exhaust, a burning facility) could dominate the spatial aggregate if only one or two readings are available. The three-sensor minimum was found empirically to eliminate most single-sensor spike artefacts without discarding a substantial fraction of valid hours.

After applying both filters, the city-wide hourly PM2.5 target is defined as the cross-sectional median of all qualifying sensor readings in that hour. The median is computed in Pandas using the expression \texttt{df\_filtered.groupby(\textquotesingle{}timestamp\textquotesingle{}).median(){[}\textquotesingle{}value\textquotesingle{}{]}}, followed by \texttt{resample(\textquotesingle{}h\textquotesingle{}).first()} to align the result to exact UTC clock hours. The median is preferred over the arithmetic mean because it is not affected by the distribution's tail: a single sensor reporting 490 μg/m³ has no influence on the median when the remaining 50 sensors report values in the range 35--45 μg/m³, whereas it would inflate the hourly mean by approximately 9 μg/m³. All 192 stations in the spatial query are fixed-location instruments; mobile platforms with GPS tracks were excluded at the query stage by specifying \texttt{entity=fixed} in the OpenAQ API call. This filter is important because mobile sensors contribute time-varying spatial samples that cannot be consistently assigned to a clock hour at a fixed grid point. The aggregation step reduces 578,993 raw sensor-hour records to 11,553 valid city-median hourly values across the training window, corresponding to a temporal coverage rate of 85.7\% of the 13,488 available meteorological hours.

All timestamps are stored and processed in Coordinated Universal Time (UTC). The OpenAQ API returns measurements in UTC; the Open-Meteo archive likewise reports UTC hourly intervals. No time-zone conversion is applied at the data-ingest stage. Calendar-based features, hour of day, day of week, day of year, are therefore computed in UTC, which is UTC+5 minus five hours relative to Almaty Local Time (ALT). At the feature-engineering stage this introduces a systematic five-hour phase offset between calendar features and local time. The heating-season indicator, which is based on calendar month and day, is unaffected because its temporal resolution is one day rather than one hour; a five-hour UTC offset does not shift a heating-season boundary date.

\subsubsection{2.3.2. Meteorological join and timestamp alignment}\label{meteorological-join-and-timestamp-alignment}

The meteorological dataset (13,488 continuous hourly rows) and the CAMS dataset (13,488 rows) are joined to the city-median PM2.5 series on the UTC timestamp. The join is implemented as an inner join, which discards PM2.5 hours that fall outside the meteorological coverage window and discards meteorological hours where no qualifying PM2.5 reading exists. Because the meteorological window is fully continuous and the PM2.5 window is approximately coextensive, the inner join retains all 11,553 PM2.5 hours without loss.

Timestamp alignment requires one non-trivial correction. OpenAQ returns hourly measurements with the timestamp at the end of the aggregation interval, for example, a record timestamped 14:00 UTC represents the hourly mean or roll-up of measurements collected between 13:00 and 14:00 UTC. The Open-Meteo archive, by contrast, follows the standard meteorological convention of reporting instantaneous values at the top of the hour, a record timestamped 14:00 UTC represents the state of the atmosphere at that moment. Joining the two series directly on the raw timestamp would align the meteorological conditions at 14:00 UTC with the PM2.5 mean covering 13:00--14:00 UTC, introducing a systematic one-hour temporal misalignment. To correct for this, an offset of -1 hour is applied to all OpenAQ timestamps before the join: the PM2.5 mean from 13:00--14:00 is reassigned the timestamp 13:00, matching the meteorological convention. After this correction, the joined dataset correctly associates the meteorological state at the beginning of an aggregation hour with the PM2.5 mean measured during that hour.

\subsubsection{2.3.3. Feature construction}\label{feature-construction}

Feature construction is implemented across three modules: \texttt{src/features/temporal.py} (calendar and cyclical encodings, heating-season indicator), \texttt{src/features/meteorological.py} (BLH transforms, wind decomposition, temperature gradient), and \texttt{src/features/pipeline.py} (autoregressive lags, rolling means, final assembly). All feature operations are applied in a fixed order after the inner join to prevent ambiguity in column availability at any intermediate step. The final feature matrix contains 27 columns, described below by feature family.

\textbf{Autoregressive PM2.5 lags.} Five lagged copies of the city-median PM2.5 series are constructed using \texttt{pandas.Series.shift(h)} for h $\in$ \{1, 3, 6, 12, 24\} hours. Shifting by h positions aligns the target at time t with the observation at time t − h. Lag features encode the recent state of the PM2.5 time series and allow the model to exploit autocorrelation at multiple timescales. The 1-hour lag captures sub-hourly concentration trends and is the strongest single predictor at the 6-hour horizon. The 24-hour lag is the most predictively important feature at short horizons for the specific diurnal dynamics of Almaty: morning combustion peaks at approximately 07:00--09:00 local time recur predictably on consecutive winter days in the absence of synoptic-scale disturbance, giving the 24-hour lag a high correlation with next-day conditions at the same local hour.

\textbf{Rolling mean features.} Three rolling means are constructed with window sizes w $\in$ \{3, 6, 24\} hours, implemented as \texttt{pm25\_series.shift(1).rolling(w).mean()}. The \texttt{shift(1)} operation before rolling ensures that the window ends at t − 1, not at t; without this shift, the rolling mean at time t would include the value y\_t itself, constituting look-ahead data leakage in the regression formulation where y\_t is also part of the target construction. The three window sizes were chosen to capture distinct temporal scales of PM2.5 dynamics in Almaty. The 3-hour window captures the near-term trend direction, whether concentrations are currently rising, stable, or falling, which is informative about the phase of an inversion episode. The 6-hour window covers the typical onset-to-peak duration of a single-night radiation inversion: a stable 6-hour mean at elevated concentrations indicates that an inversion is entrenched rather than newly forming, which carries different implications for the 12-hour and 24-hour forecast horizons. The 24-hour rolling mean captures multi-day persistence patterns characteristic of blocking anticyclones, during which PM2.5 accumulates over three to five days and remains elevated even during afternoon mixing hours.

\textbf{BLH-derived features.} Raw BLH values range from approximately 10 to 3,935 m in the training window, following a strongly right-skewed distribution with a median of approximately 620 m and a mean of approximately 850 m. Three BLH-derived features are constructed. The first is a log transform: \texttt{blh\_log\ =\ log(BLH\ +\ 1)}, implemented with \texttt{numpy.log1p}. The logarithm compresses the upper tail of the distribution and amplifies distinctions at the physically critical low end: the difference between 10 m and 50 m (a factor of five, representing the difference between a deep winter inversion and a shallow but developing one) is more physically meaningful than the difference between 2,000 m and 2,040 m. The second is the hour-over-hour delta: \texttt{blh\_delta\ =\ BLH\_t\ −\ BLH\_\{t−1\}}, computed with \texttt{pandas.Series.diff(1)}. A large negative delta (e.g., −200 m/h) indicates a rapidly collapsing boundary layer, a leading indicator of PM2.5 accumulation in the next two to six hours. A large positive delta indicates inversion breakdown and typically precedes a PM2.5 decline within three to six hours. The third feature is the raw BLH value itself, preserved alongside the transform and delta so that the model has access to the absolute mixing depth as well as its rate of change. Together these three features encode the primary physical mechanism linking meteorology to PM2.5 concentration in Almaty's winter regime.

\textbf{Wind decomposition.} Scalar wind speed and direction are decomposed into orthogonal Cartesian components using the standard meteorological convention: \texttt{u\ =\ wind\_speed\ ×\ sin(wind\_direction\_rad)} (positive eastward) and \texttt{v\ =\ wind\_speed\ ×\ cos(wind\_direction\_rad)} (positive northward). The raw direction angle is not used directly as a feature because it is a circular variable: a naive linear model treats 1° and 359° as opposite ends of a linear scale, when in fact they represent nearly identical north-north-easterly directions separated by only 2°. The (u, v) decomposition removes this discontinuity and allows the model to learn directional effects continuously. For Almaty, northerly and north-westerly winds, characterised by negative u and positive v components, advect cleaner continental air from the steppe and correlate with PM2.5 reductions of 15--30 μg/m³ compared to calm conditions. South-easterly drainage flows from the Zailiiskiy Alatau correlate with valley recirculation and PM2.5 accumulation.

\textbf{Temperature gradient.} An hour-over-hour temperature change \texttt{temp\_gradient\ =\ T\_t\ −\ T\_\{t−1\}} is computed with \texttt{pandas.Series.diff(1)}. The temperature gradient serves as a proxy for the rate of surface cooling, which drives nocturnal inversion formation through longwave radiative emission at the ground. Rapid cooling, a strongly negative gradient, indicates conditions conducive to inversion onset within the next two to four hours. In the training set, hours with a temperature gradient below −1.5°C/h are followed by PM2.5 increases averaging 12 μg/m³ over the subsequent three-hour period, which justifies including this derived feature alongside the raw temperature value.

\textbf{Cyclical calendar encodings.} Hour of day (integer range 0--23), day of week (0--6), and day of year (1--365) are each encoded as sine-cosine pairs using the formula (sin(2πx/P), cos(2πx/P)), where P is the period of the variable (24, seven, and 365, respectively). This encoding maps each integer to a unique point on the unit circle, ensuring that the encoded representation is periodic: hour 23 and hour 0 are adjacent on the circle rather than at opposite ends of a linear range. A raw integer encoding of hour of day would present the transition from 23 to 0 as a discontinuity of magnitude 23, comparable to the entire daily range, which would mislead a linear model. The sine-cosine pair also uniquely determines the position on the circle for any P, a property that the raw integer encoding inherently violates at the boundaries. The six cyclical features (two per calendar variable) replace the three raw integer columns. Day-of-week encodings are included because the distinction between weekday traffic patterns and weekend patterns is weakly associated with PM2.5 even in the heating season, though this effect is substantially weaker than the heating-season signal encoded by the binary indicator.

\textbf{Heating-season indicator.} A binary feature (\texttt{heating\_season}) is assigned the value one for all hours falling between 15 October and 31 March inclusive, corresponding to the nominal Almaty residential and district heating season when coal combustion at centralised heat plants and household boilers constitutes the dominant PM2.5 emission source. The indicator is zero for the remaining months (April through October), when vehicular emissions and secondary aerosol formation drive a lower and less variable PM2.5 background. This feature allows both the Ridge regression and XGBoost models to learn a level-shift in PM2.5 associated with the heating season independently of the meteorological predictors.

After feature construction and \texttt{dropna()} removal of the first 24 rows (which have undefined lag values due to the longest lag window of 24 hours), the feature matrix contains 11,553 rows and 27 columns. Rows with any remaining missing values, arising from sensors that were offline at specific hours and thus produced no city-median, are also discarded, bringing the usable dataset to approximately 11,200 rows before the train-test split.

\subsubsection{2.3.4. Feature selection and dimensionality}\label{feature-selection-and-dimensionality}

The 27-feature set was defined by domain knowledge and physical reasoning rather than by automated feature selection algorithms. At 27 dimensions the dataset does not present curse-of-dimensionality concerns: with approximately 9,200 training samples and 27 features, the sample-to-feature ratio is approximately 341:1, which is well above the threshold at which regularised models begin to overfit. No dimensionality reduction, principal component analysis, autoencoders, or filter-based selection, was applied at any stage.

The deliberate choice to use a domain-knowledge-driven feature set rather than an automated selection procedure has two justifications. First, it keeps the model interpretable: each of the 27 features corresponds to a physically or statistically motivated quantity (a specific lag, a physical transform, a known meteorological driver), and the XGBoost feature importance scores in §3.4 can therefore be interpreted in terms of the underlying processes rather than in terms of anonymous principal components. Second, it ensures a fair comparison between Ridge regression and XGBoost. If the two models were presented with different feature subsets, for example, if a linear-model-specific selection step removed features that are collinear under linearity assumptions but informative under tree-based splits, any performance difference between the models would be confounded with the feature-selection difference rather than reflecting the models' intrinsic capacities. By presenting both models with the identical 27-feature matrix, any difference in forecast accuracy is attributable to differences in model class (linear versus non-linear additive ensemble) rather than to differences in input information.

\subsection{2.4. Models and baselines}\label{models-and-baselines}

Four models are evaluated for each forecast horizon h $\in$ \{6, 12, 24\} hours. In each case, the target variable is y\_\{t+h\}: the city-median PM2.5 concentration h hours after the feature observation time t. The 12 model × horizon combinations are evaluated on the same 2,311-hour held-out test set using the metrics defined in §2.5.

\subsubsection{2.4.1. Persistence baseline}\label{persistence-baseline}

The persistence baseline requires no training step, no parameters, and no data beyond the most recent observation. Its forecast for any horizon h is:

ŷ\_\{t+h\} = y\_t (1)

That is, the prediction is the last observed city-median PM2.5 value, regardless of horizon. Persistence is the appropriate lower bound for any data-driven forecaster: a model that cannot outperform persistence at a given horizon provides no marginal value over a trivial no-model rule. For short horizons where PM2.5 exhibits high autocorrelation, the lag-1 autocorrelation in the training set is 0.94, persistence is a competitive baseline that many simple models fail to beat.

The performance of the persistence baseline varies markedly across horizons in Almaty's diurnal PM2.5 cycle. At the 6-hour horizon, where autocorrelation remains high, persistence typically performs well during multi-day inversion episodes but degrades during rapid inversion onset or breakup transitions. At the 12-hour horizon, persistence encounters a structural disadvantage: high morning concentrations (at 07:00--09:00 local time) are weakly anti-correlated with values 12 hours later (19:00--21:00 local time), when the boundary layer often mixes more thoroughly in winter afternoons. At the 24-hour horizon, the diurnal periodicity partially restores the autocorrelation (conditions at 08:00 today are moderately similar to conditions at 08:00 tomorrow during stable winter regimes), but multi-day trend changes mean persistence under-performs at this horizon during transition days. The anti-persistence effect at the 12-hour horizon is documented quantitatively in §3.3.

Because the persistence baseline requires zero parameters and zero training data, it is the most conservative possible lower bound. Any model that requires training data but still fails to outperform persistence should not be deployed, as it would produce worse forecasts than no model at all. This criterion is used in §3.3 to evaluate the marginal value of the CAMS external baseline.

\subsubsection{2.4.2. Ridge regression}\label{ridge-regression}

A Ridge regression model (L2-regularised linear regression) is trained on the same 27-feature matrix as the XGBoost model. Ridge regression minimises the regularised least-squares objective:

L(β) = Σᵢ (yᵢ − xᵢᵀβ)² + α × Σⱼ βⱼ² (2)

where β is the vector of regression coefficients, xᵢ is the feature vector for observation i, yᵢ is the target, and α is the regularisation strength. With α = 1.0, the L2 penalty prevents individual coefficients from growing large in response to collinear feature pairs. Several features in the 27-column matrix are highly correlated by construction: \texttt{pm25\_lag1} and \texttt{pm25\_roll3} share three overlapping observations; \texttt{blh} and \texttt{blh\_log} are a deterministic transform of one another; \texttt{u} and \texttt{v} are algebraic combinations of the same two meteorological variables. Without L2 regularisation, ordinary least squares would assign large positive and negative coefficients to collinear pairs, producing numerically unstable predictions. The Ridge penalty bounds the coefficient norm and stabilises the solution.

Because Ridge regression is not scale-invariant, all 27 features are standardised to zero mean and unit variance before fitting. Standardisation is implemented using \texttt{sklearn.preprocessing.StandardScaler} fitted exclusively on the training set (the first 9,242 rows). The resulting mean and standard deviation parameters are stored alongside the trained model coefficients and applied to the test set before prediction. Applying test-set values to training-set standardisation parameters ensures that no information from the test period leaks into the model fitting step. The regularisation parameter α = 1.0 was selected as a standard weakly informative value without cross-validated grid search; the primary purpose of the Ridge model is to serve as a linear baseline against which to assess the marginal contribution of tree-based non-linearity, rather than to maximise absolute forecast accuracy.

\subsubsection{2.4.3. CAMS external baseline}\label{cams-external-baseline}

CAMS-Global PM2.5 output at the target location and time t+h is used directly as a forecast without any training step. As described in §2.1.3, the CAMS output is a global chemistry-transport model reanalysis at approximately 40 km resolution, accessed at the grid point closest to Almaty (approximately 43.0°N, 76.5°E). Because the CAMS grid point is located approximately 30 km from the city centre, and because the 40 km horizontal resolution cannot resolve the Almaty basin topography (the city is bounded on three sides by the Zailiiskiy Alatau mountain range at elevations exceeding 3,500 m), CAMS PM2.5 represents a spatially smoothed estimate that does not capture the sub-grid inversion dynamics responsible for the most extreme PM2.5 episodes. This grid-point offset may introduce a systematic bias in CAMS forecasts for local inversion events, and the \textasciitilde30 km offset is larger than the typical horizontal scale of the inversion pool during severe episodes.

The CAMS forecast serves as an evidence-based external baseline representing the best available alternative for a user without access to local sensor data. Its inclusion allows the present study to directly measure the information gain from the 190-sensor LCS network: the difference in performance between the locally trained models and CAMS at each horizon quantifies the value added by ground-level observation in this domain.

\subsubsection{2.4.4. XGBoost}\label{xgboost}

The primary model is an XGBoost gradient-boosted tree ensemble \cite{chen2016xgboost}. Three separate models are trained, one per forecast horizon h $\in$ \{6, 12, 24\} hours, because fitting independent models per horizon allows the tree structure, depth, and feature importances to adapt to the specific statistical characteristics of each prediction problem, without the multi-output constraints that a single model covering all horizons would impose.

The gradient boosting algorithm works as follows. The ensemble is initialised with a constant prediction equal to the training-set mean of the target variable. At each of the T = 600 boosting rounds, a new regression tree is fit to the negative gradient of the loss function evaluated on the current ensemble's predictions. For mean squared error loss, the negative gradient is simply the residual (y\_i − ŷ\_i), so each tree is fit to the current prediction errors. The new tree is added to the ensemble with the learning rate η = 0.05 scaling its contribution, which shrinks each step and requires more trees to converge but generally produces better generalisation than higher learning rates. The final prediction is the sum of the constant initialisation and the weighted contributions of all 600 trees.

Each of the following hyperparameter choices was made by reasoning from the training-set size and the structure of the prediction problem rather than by automated search. The number of estimators, \texttt{n\_estimators\ =\ 600}, was selected to balance model capacity against training time: at a learning rate of 0.05, the ensemble typically requires 400--700 rounds to converge on tabular datasets of this size (approximately 9,200 training samples), and 600 provides adequate capacity while keeping training time to 4--8 seconds per model on a standard laptop CPU. Early stopping was not implemented, to keep the evaluation protocol simple and to ensure that all three per-horizon models are trained for exactly the same number of rounds. The maximum tree depth, \texttt{max\_depth\ =\ 6}, matches the XGBoost library default; depth-8 models were tested during development and showed higher test-set RMSE at the 24-hour horizon compared to depth-6 models, consistent with overfitting to the smaller within-hour autocorrelation signal at longer horizons.

The stochastic sampling parameters \texttt{subsample\ =\ 0.8} and \texttt{colsample\_bytree\ =\ 0.8} implement standard stochastic gradient boosting: each tree is trained on a random 80\% subsample of the training instances (row sampling) and a random 80\% of the features (column sampling). Row sampling reduces the correlation between successive trees and introduces ensemble diversity; column sampling provides an effect analogous to the random feature selection in random forests. The minimum leaf weight parameter \texttt{min\_child\_weight\ =\ 5} requires that any tree leaf contain at least five training samples weighted by their Hessian (second derivative of the loss, which equals one for MSE). This prevents the model from creating highly specific split paths that match only one or two training samples, a pattern observed in early experiments during the AirGradient firmware-outage period, when gaps of several hours produced small isolated clusters of low-PM2.5 values that the model would fit with leaf-level predictions far outside the normal range. The regularisation parameters \texttt{reg\_alpha\ =\ 0.1} (L1 penalty on leaf weights) and \texttt{reg\_lambda\ =\ 1.0} (L2 penalty on leaf weights) apply an elastic-net-like penalty at the leaf score level. The L2 term dominates and provides the primary regularisation; the small L1 term encourages exact-zero leaf weights for leaves that contribute negligible predictive signal, simplifying the effective tree structure.

The trained models are serialised using Python's \texttt{joblib.dump} function to files approximately 2 MB each at \texttt{models/xgboost\_h\{6,12,24\}.joblib}. The Streamlit application loads these files at startup using \texttt{joblib.load} and uses them to generate forecasts from the current feature vector without re-training.

No systematic hyperparameter search was performed. The chosen parameters represent a standard configuration for tabular regression problems of this scale, and the primary scientific objective of the study, demonstrating that a locally trained gradient-boosting model outperforms the global CTM baseline and quantifying the horizon-dependent degradation in forecast skill, does not require maximising absolute accuracy through exhaustive tuning.

\subsection{2.5. Evaluation metrics and validation strategy}\label{evaluation-metrics-and-validation-strategy}

\subsubsection{2.5.1. Evaluation metrics}\label{evaluation-metrics}

Four scalar metrics are reported for each model × horizon combination. The metrics were selected to provide complementary perspectives on forecast quality: RMSE is sensitive to large errors, MAE characterises the typical error, R² normalises performance relative to a trivial baseline, and MAPE provides a scale-free reference.

\textbf{Root mean squared error (RMSE):}

RMSE = √((1/n) × Σᵢ(yᵢ − ŷᵢ)²) (3)

RMSE weights errors quadratically, so a single error of 40 μg/m³ contributes as much to RMSE as 16 errors of 10 μg/m³ each. This quadratic emphasis is appropriate for this application because the operationally critical forecast events are inversion-onset spikes, rapid concentration increases of 50--100 μg/m³ over two to four hours, where large errors have the greatest public-health consequence. RMSE is reported in μg/m³ and is directly comparable across models and horizons. For the best-performing model (XGBoost at 6 hours), RMSE = 23.44 μg/m³ on the held-out test set.

\textbf{Mean absolute error (MAE):}

MAE = (1/n) × Σᵢ\textbar yᵢ − ŷᵢ\textbar{} (4)

MAE provides the expected absolute forecast error in μg/m³ and is less sensitive to outliers than RMSE. The ratio RMSE/MAE indicates whether forecast errors are concentrated in a small number of extreme events (RMSE \textgreater\textgreater{} MAE, indicating a heavy-tailed error distribution) or distributed more uniformly across time (RMSE ≈ MAE). A ratio substantially greater than one suggests that model failures are concentrated in specific meteorological regimes, which motivates further regime-stratified analysis in §3.5.

\textbf{Coefficient of determination (R²):}

R² = 1 − Σᵢ(yᵢ − ŷᵢ)² / Σᵢ(yᵢ − ȳ)² (5)

R² measures the proportion of the test-set variance in y explained by the model's predictions. A value of 1.0 indicates a perfect forecast; a value of 0 indicates performance equivalent to the unconditional test-set mean predictor; a negative value indicates performance worse than that trivial mean. R² of 0.5 means the model explains half the variance of PM2.5 on the test set, which is the result achieved by XGBoost at the 6-hour horizon. By contrast, R² ≤ −0.18 for all CAMS horizons means that CAMS predictions are consistently less accurate than simply predicting the test-set mean PM2.5, a result attributable to CAMS systematically underestimating winter inversion concentrations that are below its 40 km grid resolution.

\textbf{Mean absolute percentage error (MAPE):}

MAPE = (100/n) × Σᵢ\textbar yᵢ − ŷᵢ\textbar{} / max(\textbar yᵢ\textbar, ε) (6)

where ε = 1 μg/m³ is a small constant added to prevent division by zero. MAPE is reported for completeness but is a less reliable evaluation metric in this dataset. The test set includes April 2026 hours where PM2.5 values approach 10--15 μg/m³ during the spring transition; a 5 μg/m³ absolute error at these low concentrations produces a 33--50\% relative error. MAPE therefore reflects the PM2.5 concentration level as much as the model's performance, and is disproportionately influenced by the low-concentration spring period. RMSE and R² are the primary evaluation criteria referenced in the comparative analysis of Chapter 3.

\subsubsection{2.5.2. Temporal train-test split}\label{temporal-train-test-split}

The 11,553-hour feature matrix is partitioned at the 80th percentile of the temporal index: the first 9,242 hours (approximately 2024-10-02 through 2025-12-08) constitute the training set, and the remaining 2,311 hours (approximately 2025-12-09 through 2026-04-14) constitute the held-out test set. No random shuffling is applied at any stage. In time-series regression, random shuffling of the train-test partition would allow future lag features, lag-1, lag-3, lag-6, lag-12, and lag-24 PM2.5 values that encode information about the test-period target variable, to appear in the training set, a form of data leakage that would artificially inflate measured accuracy by an amount proportional to the test-period autocorrelation.

The 80/20 split ratio was chosen over the alternatives of 70/30 and 90/10. A 70/30 split would extend the test period to approximately 3,466 hours and the test start date to approximately October 2025, which would exclude from the training set the spring 2025 transition period (April--June 2025), the only examples of warm-season PM2.5 dynamics available in the dataset. Excluding these months from training would leave the model without exposure to the low-BLH, low-humidity summer conditions, increasing generalisation error during April 2026 test hours. A 90/10 split would compress the test period to approximately 1,155 hours, spanning approximately 48 days from late February through April 2026. A 48-day test window might coincidentally exclude or over-represent specific inversion episodes, producing metric estimates with higher variance than the 2,311-hour window. The 2,311-hour test period spans approximately 96 days, covering the late December 2025 through April 2026 period, which includes multiple full inversion episodes, several major PM2.5 exceedance events, and the spring transition to cleaner conditions. This range provides representative and stable estimates of all four metrics.

The test period was deliberately aligned with the most operationally demanding season: winter 2025--26, when PM2.5 concentrations are highest, variability is greatest, and the forecast challenge is hardest. A model that achieves acceptable performance on this test set provides stronger evidence of practical utility than one evaluated on a summer period where PM2.5 variability is low and persistence is trivially effective. No k-fold cross-validation was applied, because implementing k-fold validation in a time-series setting requires careful gap management between training and validation windows of each fold to prevent leakage through the lag features; the single temporal split avoids this complexity while providing a test set large enough (20\% of the dataset) to yield stable metric estimates without resampling.

\subsection{2.6. Ethical and data-governance considerations}\label{ethical-and-data-governance-considerations}

All data sources used in the present study are published under the Creative Commons Attribution 4.0 International licence (CC BY 4.0). The CC BY 4.0 terms permit use, redistribution, and the creation of derivative works, including trained models, transformed datasets, and deployed applications, provided that the original sources are attributed. The three primary sources (OpenAQ \cite{openaq_platform}, Open-Meteo \cite{zippenfenig_openmeteo_2023}, and CAMS via Open-Meteo) each satisfy this condition, and attribution is provided in the bibliography and in the application interface. No data source requires a formal data-sharing agreement beyond registration of a free API key (OpenAQ) or carries a commercial-use restriction applicable to academic research (Open-Meteo, CAMS).

No personal data are collected or processed at any stage of the pipeline. The LCS sensors in the OpenAQ dataset are fixed-location, publicly cited instruments; the present study uses only their measurement values and spatial coordinates, not any identifying information about operators or owners. The city-median aggregation step (§2.3.1) discards all per-sensor identifying information before the modelling stage: the feature matrix contains the computed city-median value and no other sensor-level attribute. The Streamlit application does not require user authentication, does not collect usage metrics, and does not transmit user data to any third party. No access to private premises, individual persons, or sensitive infrastructure is required by the data collection protocol.

The use of uncalibrated LCS data raises a consideration that extends beyond licence terms. Low-cost optical particle counters of the AirGradient and Clarity type are not traceable to a gravimetric mass standard (the NIST PM2.5 reference method) and may under- or over-report PM2.5 by factors of 0.5 to 2.0 depending on aerosol composition, relative humidity, and instrument aging \cite{who_aqg_2021}. Forecasts produced by a model trained on these observations inherit this uncertainty: predicted concentrations are indicative rather than reference-grade. For this reason, the forecasting system described here is not intended for use as the sole basis for regulatory air quality decisions, emergency alerts, or health advisories that carry legal or medical weight. The Streamlit application interface incorporates a visible disclaimer, present in the application's sidebar and on the forecast card, stating that PM2.5 values are derived from low-cost sensors and are indicative only, consistent with the guidance in §3.5.3 of the application design chapter.

The choice to publish the pipeline code under an open-source licence and to rely exclusively on openly licensed input data is a reproducibility decision as much as an ethical one. Any researcher with valid API credentials and the repository code can rebuild the full dataset from scratch, retrain all models, and independently verify the results reported in Chapter 3. This commitment to full computational reproducibility is particularly important given that the dataset cannot be packaged in the repository due to its size (the SQLite database is approximately 50 MB). The gitignore configuration ensures that raw Parquet files and the SQLite database are never accidentally committed to version control, keeping the repository lightweight and preventing inadvertent data redistribution that could conflict with sensor operators' terms of service.
